\documentclass[11pt,a4paper]{article}
\usepackage[utf8]{inputenc}
\usepackage[T1]{fontenc}
\usepackage[italian]{babel}
\usepackage{lmodern}
\usepackage{microtype}
\usepackage[margin=2cm,headheight=14pt]{geometry}
\usepackage{amsmath,amssymb,amsthm,mathtools,amsfonts,bm,siunitx}
\usepackage{booktabs,tabularx,enumitem,float}
\usepackage{xcolor}
\usepackage[most]{tcolorbox}
\usepackage{fancyhdr}
\usepackage{listings}
\usepackage{hyperref}
\usepackage{bookmark}

\definecolor{navyblue}{RGB}{10, 45, 90}
\definecolor{coolblue}{RGB}{28, 110, 164}
\definecolor{forestgreen}{RGB}{20, 115, 60}
\definecolor{warmamber}{RGB}{195, 110, 10}
\definecolor{crimsonred}{RGB}{175, 30, 45}
\definecolor{subtlegray}{RGB}{245, 247, 250}
\definecolor{codebg}{RGB}{240, 242, 245}

\hypersetup{
    colorlinks=true,
    linkcolor=navyblue,
    urlcolor=coolblue,
    citecolor=forestgreen
}

\pagestyle{fancy}
\fancyhf{}
\fancyhead[L]{\textbf{\textsf{\textcolor{navyblue}{Statistica I --- Soluzione Ufficiale Dettagliata}}}}
\fancyhead[R]{\textsf{\textcolor{coolblue}{III Appello (Luglio 2020)}}}
\fancyfoot[C]{\thepage}
\renewcommand{\headrulewidth}{0.6pt}

\newtcolorbox{problemabox}[1]{
    enhanced,
    breakable,
    colback=white,
    colframe=navyblue,
    coltitle=white,
    fonttitle=\bfseries\sffamily,
    title={#1},
    attach boxed title to top left={yshift=-2mm, xshift=3mm},
    boxed title style={colback=navyblue, boxrule=0pt, sharp corners, rounded corners=north},
    boxrule=1.2pt,
    sharp corners,
    rounded corners=south,
    top=4mm, bottom=3mm, left=3mm, right=3mm
}

\newtcolorbox{soluzionebox}[1][Svolgimento Dettagliato]{
    enhanced,
    breakable,
    colback=subtlegray,
    colframe=forestgreen!80!black,
    coltitle=white,
    fonttitle=\bfseries\sffamily,
    title={#1},
    attach boxed title to top left={yshift=-2mm, xshift=3mm},
    boxed title style={colback=forestgreen!80!black, boxrule=0pt, sharp corners, rounded corners=north},
    boxrule=1pt,
    sharp corners,
    rounded corners=south,
    top=4mm, bottom=3mm, left=3mm, right=3mm
}

\newtcolorbox{esamebox}[1][Consiglio per l'Esame \& Aspetti Chiave]{
    enhanced,
    breakable,
    colback=crimsonred!4!white,
    colframe=crimsonred,
    coltitle=crimsonred,
    fonttitle=\bfseries\sffamily,
    title={#1},
    attach boxed title to top left={yshift=-1.5mm, xshift=2mm},
    boxed title style={colback=white, boxrule=0.8pt, colframe=crimsonred},
    boxrule=0.8pt,
    leftrule=3.5pt,
    sharp corners,
    top=3.5mm, bottom=2.5mm, left=3mm, right=3mm
}

\lstset{
    backgroundcolor=\color{codebg},
    basicstyle=\ttfamily\small,
    breaklines=true,
    frame=single,
    rulecolor=\color{navyblue!40},
    keywordstyle=\color{navyblue}\bfseries,
    commentstyle=\color{forestgreen}\itshape,
    stringstyle=\color{warmamber},
    showstringspaces=false
}

\newcommand{\EE}{\mathbb{E}}
\newcommand{\Prob}{\mathbb{P}}
\newcommand{\Var}{\operatorname{Var}}
\newcommand{\Cov}{\operatorname{Cov}}
\newcommand{\dx}{\mathrm{d}x}

\begin{document}

\begin{center}
    {\LARGE\bfseries\color{navyblue} Università di Pisa --- Corso di Laurea in Ingegneria Gestionale}\\[0.4em]
    {\Large\bfseries Statistica I (A.A. 2019/2020) --- III Appello (Luglio 2020)}\\[0.3em]
    {\large\textsf{Soluzione Completa, Rigorosa e Commentata Passo-Passo}}
\end{center}
\vspace{0.5em}
\hrule height 1.2pt
\vspace{1.5em}

% ==============================================================================
% PROBLEMA 1
% ==============================================================================
\begin{problemabox}{Problema 1 (Testo Integrale)}
Siano $X,Y$ due variabili aleatorie standardizzate e indipendenti.
\begin{enumerate}[label=(\roman*)]
    \item Si calcolino il valore atteso e la varianza di $Z=X+3Y$ e di $U=XY$.
    \item Si dimostri che vale la disuguaglianza $\Prob(|X+Y| \ge 4) \le \frac{1}{8}$.
    \item Siano $U,V$ indipendenti e di Bernoulli di parametro rispettivamente $\frac{1}{2}$ e $\frac{1}{3}$. Si calcoli la probabilità $\Prob(X=2, Y=-\sqrt{3/2})$ nel caso in cui $X$ e $Y$ siano state ottenute standardizzando $U$ e $V$.
\end{enumerate}
\end{problemabox}

\begin{soluzionebox}
\textbf{1. Proprietà di variabili standardizzate indipendenti:}
Essendo $X$ e $Y$ standardizzate, abbiamo per definizione:
\[
\EE[X] = 0, \quad \EE[Y] = 0, \quad \Var(X) = 1, \quad \Var(Y) = 1 \implies \EE[X^2] = 1, \quad \EE[Y^2] = 1.
\]
L'indipendenza stocastica implica $\Cov(X,Y) = 0$ e $\EE[g(X)h(Y)] = \EE[g(X)]\EE[h(Y)]$.

\textbf{2. Risoluzione dei quesiti:}
\begin{enumerate}[label=(\roman*)]
    \item \textbf{Valore atteso e varianza di $Z = X+3Y$ e $U = XY$:}
    \begin{itemize}
        \item Per $Z = X + 3Y$:
        \[
        \EE[Z] = \EE[X] + 3\EE[Y] = 0 + 3(0) = \bm{0},
        \]
        \[
        \Var(Z) = \Var(X) + 3^2 \Var(Y) = 1 + 9(1) = \bm{10}.
        \]
        \item Per $U = XY$:
        \[
        \EE[U] = \EE[XY] = \EE[X]\EE[Y] = 0 \cdot 0 = \bm{0},
        \]
        \[
        \Var(U) = \EE[U^2] - (\EE[U])^2 = \EE[X^2 Y^2] - 0 = \EE[X^2]\EE[Y^2] = 1 \cdot 1 = \bm{1}.
        \]
    \end{itemize}

    \item \textbf{Dimostrazione tramite Disuguaglianza di Chebyshev:}
    Definiamo la variabile aleatoria $W = X + Y$.
    Calcoliamo media e varianza di $W$:
    \[
    \EE[W] = \EE[X] + \EE[Y] = 0, \qquad \Var(W) = \Var(X) + \Var(Y) = 1 + 1 = 2.
    \]
    La disuguaglianza di Chebyshev stabilisce che per ogni costante $k > 0$:
    \[
    \Prob(|W - \EE[W]| \ge k) \le \frac{\Var(W)}{k^2}.
    \]
    Applicando la disuguaglianza con $k = 4$ e $\EE[W]=0$:
    \[
    \Prob(|X + Y| \ge 4) \le \frac{\Var(W)}{4^2} = \frac{2}{16} = \bm{\frac{1}{8}}. \qquad \blacksquare
    \]

    \item \textbf{Calcolo di $\Prob(X=2, Y=-\sqrt{3/2})$:}
    Consideriamo $U \sim \operatorname{Bernoulli}(1/2)$ e $V \sim \operatorname{Bernoulli}(1/3)$:
    \begin{itemize}
        \item $\EE[U] = 1/2$, $\Var(U) = 1/4 \implies \sigma_U = 1/2$.
        La variabile standardizzata è $X = \frac{U - 1/2}{1/2} = 2U - 1$.
        Poiché $U \in \{0, 1\}$, la variabile $X$ può assumere unicamente i valori $\{2(0)-1, 2(1)-1\} = \{-1, +1\}$.
        Il supporto di $X$ è $\operatorname{supp}(X) = \{-1, 1\}$.
        \item $\EE[V] = 1/3$, $\Var(V) = 2/9 \implies \sigma_V = \frac{\sqrt{2}}{3}$.
        $Y = \frac{V - 1/3}{\sqrt{2}/3} \in \left\{ -\frac{1}{\sqrt{2}}, \sqrt{2} \right\}$.
    \end{itemize}
    Dato che il valore $2 \notin \operatorname{supp}(X)$, l'evento $\{X = 2\}$ è l'evento impossibile ($\Prob(X=2) = 0$).
    Pertanto:
    \[
    \Prob\left(X=2, \, Y=-\sqrt{\frac{3}{2}}\right) \le \Prob(X = 2) = \bm{0}.
    \]
\end{enumerate}
\end{soluzionebox}

% ==============================================================================
% PROBLEMA 2
% ==============================================================================
\begin{problemabox}{Problema 2 (Testo Integrale)}
Implementando in R il Metodo Monte Carlo si calcolino approssimativamente:
\begin{enumerate}[label=(\roman*)]
    \item l'integrale $\int_0^{1/2} \frac{e^x - (\sin(x))^2}{\sqrt{x}} \, \dx$,
    \item la serie $\sum_{x=0}^\infty \frac{2^x \sin(x)}{x!}$.
\end{enumerate}
\end{problemabox}

\begin{soluzionebox}
\begin{enumerate}[label=(\roman*)]
    \item \textbf{Integrazione Monte Carlo:}
    Riscriviamo l'integrale $I = \int_0^{1/2} g(x) \, \dx$ come valore atteso rispetto ad una v.a. $X \sim \operatorname{Unif}(0, 1/2)$:
    \[
    I = \frac{1}{2} \int_0^{1/2} g(x) \frac{1}{1/2} \, \dx = \frac{1}{2} \EE[g(X)], \qquad g(x) = \frac{e^x - \sin^2(x)}{\sqrt{x}}.
    \]
    Valore esatto dell'integrale: $I \approx \bm{1.6224}$.

    \item \textbf{Calcolo Monte Carlo della Serie infinita:}
    Moltiplichiamo e dividiamo la serie per $e^2$:
    \[
    S = \sum_{x=0}^\infty \frac{2^x \sin(x)}{x!} = e^2 \sum_{x=0}^\infty \sin(x) \frac{e^{-2} 2^x}{x!} = e^2 \, \EE[\sin(K)], \qquad K \sim \operatorname{Poiss}(\lambda=2).
    \]
    Valore analitico esatto della serie:
    \[
    S = \operatorname{Im}\left( \sum_{x=0}^\infty \frac{(2e^i)^x}{x!} \right) = \operatorname{Im}\left( e^{2 e^i} \right) = e^{2\cos(1)} \sin(2\sin(1)) \approx \bm{2.9280}.
    \]
\end{enumerate}

\textbf{Script R completo:}
\begin{lstlisting}[language=R]
set.seed(42)
N <- 1000000

# (i) Integrale Monte Carlo
x_unif <- runif(N, min = 0, max = 0.5)
g_x <- (exp(x_unif) - sin(x_unif)^2) / sqrt(x_unif)
I_hat <- 0.5 * mean(g_x)
cat(sprintf("(i) Stima Integrale: %.4f\n", I_hat)) # 1.6224

# (ii) Serie Monte Carlo tramite Poisson(2)
k_poiss <- rpois(N, lambda = 2)
S_hat <- exp(2) * mean(sin(k_poiss))
cat(sprintf("(ii) Stima Serie: %.4f\n", S_hat)) # 2.9280
\end{lstlisting}
\end{soluzionebox}

% ==============================================================================
% PROBLEMA 3
% ==============================================================================
\begin{problemabox}{Problema 3 (Testo Integrale)}
Si consideri una popolazione data dalla funzione di ripartizione
\[
F_\theta(x) = \begin{cases}
1 - e^{-2\theta\sqrt{x}} & \text{se } x > 0, \\
0 & \text{se } x \le 0,
\end{cases}
\]
con $\theta > 0$ parametro incognito. Indichiamo con $X_1, \dots, X_n$ un campione casuale estratto dalla popolazione.
\begin{enumerate}[label=(\roman*)]
    \item Si determini la funzione di densità corrispondente a $F_\theta$ e la funzione di verosimiglianza.
    \item Si propongano due stimatori puntuali $T_1$ e $T_2$ per il parametro $\theta$, usando rispettivamente il metodo di massima verosimiglianza e il metodo dei momenti.
    \item Si utilizzi il metodo dell'inversione per scrivere $X$ come funzione di $U \sim \operatorname{Unif}(0,1)$, si generi in R un campione con $\theta=1/2, n=500$ e si riporti la varianza campionaria.
\end{enumerate}
\end{problemabox}

\begin{soluzionebox}
\begin{enumerate}[label=(\roman*)]
    \item \textbf{Funzione di Densità e Verosimiglianza:}
    Derivando la CDF per $x > 0$:
    \[
    f_\theta(x) = \frac{\mathrm{d}}{\dx} F_\theta(x) = -e^{-2\theta\sqrt{x}} \left( -2\theta \cdot \frac{1}{2\sqrt{x}} \right) = \mathbf{\frac{\theta}{\sqrt{x}} e^{-2\theta\sqrt{x}} \mathbf{1}_{(0, \infty)}(x)}.
    \]
    La funzione di verosimiglianza per un campione i.i.d. $x_1, \dots, x_n$ è:
    \[
    L(\theta) = \prod_{i=1}^n f_\theta(x_i) = \theta^n \left( \prod_{i=1}^n x_i^{-1/2} \right) \exp\left( -2\theta \sum_{i=1}^n \sqrt{x_i} \right).
    \]

    \item \textbf{Stimatori puntuali $T_1$ (MLE) e $T_2$ (Metodo dei Momenti):}
    \begin{itemize}
        \item \emph{Massima Verosimiglianza ($T_1$):}
        Calcoliamo la log-verosimiglianza:
        \[
        \ell(\theta) = \log L(\theta) = n \log \theta - \frac{1}{2}\sum_{i=1}^n \log x_i - 2\theta \sum_{i=1}^n \sqrt{x_i}.
        \]
        Derivando rispetto a $\theta$ e ponendo a zero:
        \[
        \frac{\partial \ell}{\partial \theta} = \frac{n}{\theta} - 2\sum_{i=1}^n \sqrt{x_i} = 0 \implies \bm{T_1 = \widehat{\theta}_{\text{MLE}} = \frac{n}{2\sum_{i=1}^n \sqrt{X_i}} = \frac{1}{2 \overline{\sqrt{X}}}}.
        \]
        \item \emph{Metodo dei Momenti ($T_2$):}
        Operiamo il cambio di variabile $Y = 2\theta\sqrt{X} \implies X = \frac{Y^2}{4\theta^2}$.
        La CDF di $Y$ è $\Prob(Y \le y) = \Prob(X \le y^2/(4\theta^2)) = 1 - e^{-y}$, quindi $Y \sim \operatorname{Exp}(1)$.
        Ne consegue che $\EE[Y^2] = \Var(Y) + (\EE[Y])^2 = 1 + 1 = 2$.
        \[
        \EE[X] = \EE\left[ \frac{Y^2}{4\theta^2} \right] = \frac{\EE[Y^2]}{4\theta^2} = \frac{2}{4\theta^2} = \frac{1}{2\theta^2}.
        \]
        Ponendo $\EE[X] = \overline{X}_n \implies \frac{1}{2\theta^2} = \overline{X}_n \implies \bm{T_2 = \widehat{\theta}_{\text{MM}} = \frac{1}{\sqrt{2\overline{X}_n}}}$.
    \end{itemize}

    \item \textbf{Metodo dell'Inversione e Simulazione in R:}
    Poniamo $F_\theta(X) = U \sim \operatorname{Unif}(0, 1)$:
    \[
    1 - e^{-2\theta\sqrt{X}} = U \implies e^{-2\theta\sqrt{X}} = 1 - U \implies -2\theta\sqrt{X} = \log(1-U) \implies \bm{X = \frac{(-\log(1-U))^2}{4\theta^2}}.
    \]
    Poiché $1-U \sim \operatorname{Unif}(0,1)$, possiamo scrivere equivalentemente $X = \frac{(-\log U)^2}{4\theta^2}$.
\end{enumerate}

\begin{lstlisting}[language=R]
set.seed(42)
n <- 500
theta <- 0.5

# Campionamento tramite metodo dell'inversione
u <- runif(n)
x_campione <- (log(u))^2 / (4 * theta^2)

# Varianza campionaria
var_camp <- var(x_campione)
cat(sprintf("Varianza campionaria ottenuta (n=500, theta=0.5): %.4f\n", var_camp))
\end{lstlisting}
\end{soluzionebox}

% ==============================================================================
% PROBLEMA 4
% ==============================================================================
\begin{problemabox}{Problema 4 (Testo Integrale)}
Si consideri l'esperimento aleatorio in cui per $50$ volte viene estratto a caso un numero reale dall'intervallo $-1 \le x \le 1$ e poi vengono sommati tutti i numeri estratti.
\begin{enumerate}[label=(\roman*)]
    \item Si calcoli approssimativamente la probabilità che il risultato finale sia compreso fra $-\frac{1}{10}$ e $\frac{1}{5}$.
    \item Si calcoli la probabilità che non venga mai estratto un numero maggiore di $0.9$.
\end{enumerate}
\end{problemabox}

\begin{soluzionebox}
\textbf{1. Modello probabilistico:}
Siano $X_1, \dots, X_{50} \overset{\text{i.i.d.}}{\sim} \operatorname{Unif}(-1, 1)$.
I parametri del singolo elemento sono:
\[
\EE[X_i] = 0, \qquad \Var(X_i) = \frac{(1 - (-1))^2}{12} = \frac{4}{12} = \frac{1}{3}.
\]
Per la somma $S_{50} = \sum_{i=1}^{50} X_i$:
\[
\EE[S_{50}] = 50(0) = 0, \qquad \Var(S_{50}) = 50 \cdot \frac{1}{3} = \frac{50}{3} \approx 16.6667 \implies \sigma_{S_{50}} = \sqrt{\frac{50}{3}} \approx 4.0825.
\]

\textbf{2. Risoluzione dei quesiti:}
\begin{enumerate}[label=(\roman*)]
    \item \textbf{Approssimazione normale della somma (TLC):}
    Per il Teorema del Limite Centrale, $S_{50} \approx \mathcal{N}(0, 50/3)$. Standardizzando:
    \begin{align*}
    \Prob\left(-\frac{1}{10} \le S_{50} \le \frac{1}{5}\right) &= \Prob\left(\frac{-0.1 - 0}{4.0825} \le Z \le \frac{0.2 - 0}{4.0825}\right) \\
    &= \Prob(-0.0245 \le Z \le 0.0490) = \Phi(0.0490) - \Phi(-0.0245) \\
    &= \Phi(0.0490) - (1 - \Phi(0.0245)) = 0.5195 - (1 - 0.5098) = \bm{0.0293} \quad (\bm{2.93\%}).
    \end{align*}

    \item \textbf{Probabilità che nessun valore superi $0.9$:}
    Per una singola estrazione uniforme su $[-1, 1]$:
    \[
    \Prob(X_i \le 0.9) = \frac{0.9 - (-1)}{1 - (-1)} = \frac{1.9}{2} = 0.95.
    \]
    Per l'indipendenza delle 50 estrazioni:
    \[
    \Prob\left(\max_{1 \le i \le 50} X_i \le 0.9\right) = \left( \Prob(X_1 \le 0.9) \right)^{50} = (0.95)^{50} \approx \bm{0.0769} \quad (\bm{7.69\%}).
    \]
\end{enumerate}
\end{soluzionebox}

% ==============================================================================
% PROBLEMA 5
% ==============================================================================
\begin{problemabox}{Problema 5 (Testo Integrale)}
Si vuole determinare la lunghezza media $\mu$ (in cm) delle mine per matita. La misurazione di un campione di 10 mine ha fornito i dati del file csv.
\begin{enumerate}[label=(\roman*)]
    \item Si calcolino media e varianza campionaria delle misure ottenute utilizzando R.
    \item Si determini un intervallo di confidenza di livello $0.95$ per $\mu$ basandosi sul campione di dati.
    \item La ditta di produzione dichiara una lunghezza media di $13\text{ cm}$. Si decida tramite un test al livello del $5\%$ se le misurazioni effettuate sul campione supportano tale affermazione o meno.
\end{enumerate}
\end{problemabox}

\begin{soluzionebox}
\textbf{1. Procedura formale di inferenza ($t$-Student a 9 g.d.l.):}
Assumendo che le lunghezze seguano una distribuzione normale $\mathcal{N}(\mu, \sigma^2)$ con varianza non nota, si utilizzano la media campionaria $\overline{X}_{10}$ e la varianza campionaria corretta $S_{10}^2$.

\textbf{2. Codice R per l'analisi completa e la verifica delle ipotesi:}
\begin{lstlisting}[language=R]
# 1. Caricamento dati e statistiche descrittive
dati_mine <- read.csv("dati_mine.csv")
x <- dati_mine[, 1]
n <- length(x) # n = 10
x_bar <- mean(x)
s2 <- var(x)
s <- sd(x)

cat(sprintf("Media campionaria: %.4f cm\n", x_bar))
cat(sprintf("Varianza campionaria: %.4f cm^2\n", s2))

# 2. Intervallo di Confidenza al 95% per mu
alpha <- 0.05
t_crit <- qt(1 - alpha/2, df = n - 1) # t_0.025, 9 = 2.2622
se_mean <- s / sqrt(n)
ic <- c(x_bar - t_crit * se_mean, x_bar + t_crit * se_mean)
cat(sprintf("IC 95%%: [%.4f, %.4f]\n", ic[1], ic[2]))

# 3. Test Bilaterale H0: mu = 13 vs H1: mu != 13 al livello 5%
mu_0 <- 13
t_stat <- (x_bar - mu_0) / se_mean
p_value <- 2 * (1 - pt(abs(t_stat), df = n - 1))

cat(sprintf("Statistica test t: %.4f, Valore Critico: +/-%.4f\n", t_stat, t_crit))
cat(sprintf("p-value: %.4f\n", p_value))
if (abs(t_stat) > t_crit) {
  cat("Conclusione: Rifiutiamo H0. Le misurazioni NON supportano l'affermazione della ditta.\n")
} else {
  cat("Conclusione: Non rifiutiamo H0. I dati sono compatibili con una media di 13 cm.\n")
}
\end{lstlisting}
\end{soluzionebox}

\end{document}
